% Bench layout: the shield plugs on, nothing is wired.
\begin{tikzpicture}[font=\sffamily\scriptsize]
\node[board, minimum width=48mm, minimum height=26mm] (nuc) at (0,0) {};
\node[text=white, font=\sffamily\bfseries\small] at (0.2,-0.85) {NUCLEO-H7A3ZI-Q};

\node[hat, minimum width=40mm, minimum height=16mm] (sh) at (0.25,0.75) {};
\node[text=white, font=\sffamily\bfseries\small] at (0.25,1.0) {X-NUCLEO-IKS4A1};
\node[text=white, font=\tiny] at (0.25,0.6) {plugs onto the Arduino header};

\node[pinrow, minimum width=36mm] at (0.25,0.0) {};
\foreach \x in {-1.55,-1.35,...,2.05}{\fill[yellow!80!black] (\x,0.0) circle (0.35mm);}

\node[draw=black!70, fill=gray!60, minimum width=4mm, minimum height=6mm] (usb) at (-2.55,-0.45) {};
\node[ext, left=32mm of nuc, text width=26mm, minimum height=16mm] (pc)
  {Host PC\\{\tiny collects the histogram,\\computes the percentiles}};
\draw[usbcable] (usb.west) to[out=180,in=0]
  node[font=\tiny, fill=white, inner sep=1.5pt, pos=0.5, above]{one micro USB lead}
  node[font=\tiny, fill=white, inner sep=1.5pt, pos=0.5, below]{flash, then results back}
  (pc.east);

\node[umlnote, text width=34mm, anchor=west] at (3.0,0.6)
  {One shield at a time. Two would collide on addresses and on the 3V3
   budget.};

\node[umlnote, text width=62mm, below=10mm of nuc]
  {The release source for the million-event run is a timer on the board itself,
   not the sensor, so the release time is known rather than inferred. The
   sensor-driven run follows and is the real workload.};
\end{tikzpicture}
